\section{Controlled inherited-state populations}
\label{sec:model}

Let $\{1,\ldots,m\}$ be a finite type space and $Z_t\in\N^m$ the count vector.
A cell of type $i$ changes type according to a conservative generator $Q$,
dies at rate $d_i$, and is replaced at rate $b_i$ by an offspring vector with
probability generating function $F_i(x)$. Offspring types may be correlated; the
joint offspring law is part of the data. We assume a deterministic upper bound
$K$ on offspring count, and that all per-cell rates are bounded over the
admitted action set. Define
\begin{equation}\label{eq:field}
 \Phi(x)=Qx+d\odot(1-x)+b\odot(F(x)-x),
\end{equation}
where $\odot$ is the entrywise product. For binary offspring write
$F_i(x)=D_i(x,x)$ with $D_i$ the bilinear form of the ordered daughter law; the
mother's removal accounts for the term $-b_ix_i$.

Under deterministic rates, conditional independence of future offspring clocks
given their birth states gives the backward equation $\dot u=\Phi(u)$. With
terminal datum zero its solution is extinction by the horizon. A terminal datum
$q^{\mathrm{off}}$ instead represents eventual extinction under a specified
continuation after the course. For piecewise constant schedules the autonomous
flows act on this terminal datum in reverse chronological order.

\subsection{Actions, feedback and budget}

During intervention use $Q_v=Q_{\mathrm{base}}+vR$ with $R\one=0$ and
$0\le v_t\le v_{\max}$, each $Q_v$ a valid generator. We write $\PhiB$ for
\eqref{eq:field} at $v=0$ and $\Phi_v$ for the field with $Q_v$ in place of $Q$.
The scalar common action is predictable with respect to the entire population
history and any independent controller randomization. It obeys the pathwise
constraint
\begin{equation}\label{eq:budget}
 C_t=\int_0^tv_s\,ds\le B \qquad(0\le t\le T)
 \qquad\text{almost surely.}
\end{equation}
The constraint is pathwise, not an expectation constraint: replacing
\eqref{eq:budget} by $\E C_T\le B$ requires a different argument. Nor is $B$ a
sum of cell-specific exposures; the common-control budget is a single number
spent by a single actuator. The background death and reproduction rates are
fixed unless an extension explicitly changes them. This full-information
feedback class contains any controller using less information. Families need not
be independent under it.

The process is constructed by predictable thinning of bounded Poisson proposals
on a countable genealogical tree. Birth count is dominated by a
bounded-offspring pure-birth process, with expected size at most
$|Z_0|e^{b_{\max}(K-1)t}$ for $K\ge1$. Expected proposal counts on finite
intervals are finite, which proves non-explosion by localization. Empty
populations remain empty, so $0$ is absorbing and there is no immigration.
Feedback changes conditional intensities; it does not invalidate this
construction. A \emph{continuation} is a specified uncontrolled branching
source, started with fresh randomness conditional on the population at
withdrawal.

\subsection{The population generator identity}

For $x\in(0,1]^m$ put $V_x(z)=\prod_ix_i^{z_i}$. A type change, a death and a
reproduction multiply $V_x$ by $x_j/x_i$, by $1/x_i$ and by the offspring
monomial divided by $x_i$, respectively. Summing conditional rates, the actual
population generator obeys
\begin{equation}\label{eq:generator}
 \mathcal G_vV_x(z)=V_x(z)\sum_iz_i\frac{\Phi_v(x)_i}{x_i}.
\end{equation}
This is an identity on the population state space. It remains valid under
feedback because it is conditional on the current predictable action, and it
makes no assertion that cells exposed to a common controller have independent
histories; it does not factor founder probabilities. Conditional independence
is used later only for a fixed baseline continuation, once the population at
withdrawal is given.

\subsection{First moments}
\label{sec:firstmoment}

Write $L$ for the expected offspring matrix, $L_{ij}=\E[\#\{\text{daughters of
type }j\}\mid\text{mother of type }i]$, so that $L=2\Lambda$ for binary division
with single-daughter marginal $\Lambda$. The first-moment matrix at constant
action $v$ is
\begin{equation}\label{eq:meanmatrix}
 A_v=Q_v+\diag(b)\,(L-I)-\diag(d),
\end{equation}
acting on weight vectors from the left index, so that $W(z)=\sum_iz_iw_i$ has
drift $\sum_iz_i(A_vw)_i$. Equivalently, $A_v$ is the linearization of $\Phi_v$
at $x=\one$ up to sign: $D\Phi_v(\one)=-A_v$. It is a Metzler matrix. The
one-daughter marginal determines $A_v$ but not the extinction probabilities;
sibling correlation lives in the joint law $D$ and is retained throughout.

\subsection{Endpoints that must stay distinct}
\label{sec:endpoints}

Write
\[
 s_\pi(T)=\Prob_\pi(|Z_T|>0),\qquad
 r_\pi(T)=\Prob_\pi(\text{non-extinction under the specified continuation
 after }T).
\]
Then $s_\pi(T)\ge r_\pi(T)$, because extinction is absorbing and the
continuation has no immigration. A finite-time survival \emph{upper} bound
therefore also bounds eventual survival; a lower bound on eventual survival also
lower-bounds finite-time survival, provided its continuation assumptions hold.
Our principal continuation withdraws the additional eraser, restoring baseline
erasure while retaining the background demographic rates. Withdrawing the
background killing as well is a different experiment, and
Section~\ref{sec:withdrawal} keeps the two separate.

\subsection{Units}
\label{sec:units}

Reaction, division and death rates have units of inverse time. The integrated
exposure $B=\int v\,dt$ is dimensionless when $v$ is an additional first-order
rate. Neither $B$ nor $v$ is a drug mass or a concentration. A
nondimensionalization choosing a reference division rate $b_0$, setting
$\hat t=b_0t$ and dividing all rates by $b_0$ leaves $\int v\,dt$ unchanged, so
exposure statements transfer between time units; changing the clock alone does
not calibrate the relative biology. For the synthetic source of
Section~\ref{sec:source}, $b_0=.1$ and time $100$ spans ten mean division
waiting times.
